\documentclass[conference]{IEEEtran}

\usepackage{amsmath,amssymb,amsfonts}
\usepackage{algorithmic}
\usepackage{algorithm}
\usepackage{graphicx}
\usepackage{textcomp}
\usepackage{xcolor}
\usepackage{booktabs}
\usepackage{multirow}
\usepackage{hyperref}
\usepackage{cite}
\usepackage{array}
\usepackage{fancyvrb}
\usepackage{longtable}

% Pandoc tightlist support
\providecommand{\tightlist}{\setlength{\itemsep}{0pt}\setlength{\parskip}{0pt}}

% Pandoc syntax highlighting (if code blocks are used)
\newcommand{\VerbBar}{|}
\newcommand{\VERB}{\Verb[commandchars=\\\{\}]}
\DefineVerbatimEnvironment{Highlighting}{Verbatim}{commandchars=\\\{\},fontsize=\small}
\newenvironment{Shaded}{\vspace{4pt}}{\vspace{4pt}}
\newcommand{\AlertTok}[1]{\textcolor[rgb]{1.00,0.00,0.00}{\textbf{#1}}}
\newcommand{\AnnotationTok}[1]{\textcolor[rgb]{0.38,0.63,0.69}{\textbf{\textit{#1}}}}
\newcommand{\AttributeTok}[1]{\textcolor[rgb]{0.49,0.56,0.16}{#1}}
\newcommand{\BaseNTok}[1]{\textcolor[rgb]{0.25,0.63,0.44}{#1}}
\newcommand{\BuiltInTok}[1]{\textcolor[rgb]{0.00,0.50,0.00}{#1}}
\newcommand{\CharTok}[1]{\textcolor[rgb]{0.25,0.44,0.63}{#1}}
\newcommand{\CommentTok}[1]{\textcolor[rgb]{0.38,0.63,0.69}{\textit{#1}}}
\newcommand{\CommentVarTok}[1]{\textcolor[rgb]{0.38,0.63,0.69}{\textbf{\textit{#1}}}}
\newcommand{\ConstantTok}[1]{\textcolor[rgb]{0.53,0.00,0.00}{#1}}
\newcommand{\ControlFlowTok}[1]{\textcolor[rgb]{0.00,0.44,0.13}{\textbf{#1}}}
\newcommand{\DataTypeTok}[1]{\textcolor[rgb]{0.56,0.13,0.00}{#1}}
\newcommand{\DecValTok}[1]{\textcolor[rgb]{0.25,0.63,0.44}{#1}}
\newcommand{\DocumentationTok}[1]{\textcolor[rgb]{0.73,0.13,0.13}{\textit{#1}}}
\newcommand{\ErrorTok}[1]{\textcolor[rgb]{1.00,0.00,0.00}{\textbf{#1}}}
\newcommand{\ExtensionTok}[1]{#1}
\newcommand{\FloatTok}[1]{\textcolor[rgb]{0.25,0.63,0.44}{#1}}
\newcommand{\FunctionTok}[1]{\textcolor[rgb]{0.02,0.16,0.49}{#1}}
\newcommand{\ImportTok}[1]{\textcolor[rgb]{0.00,0.50,0.00}{\textbf{#1}}}
\newcommand{\InformationTok}[1]{\textcolor[rgb]{0.38,0.63,0.69}{\textbf{\textit{#1}}}}
\newcommand{\KeywordTok}[1]{\textcolor[rgb]{0.00,0.44,0.13}{\textbf{#1}}}
\newcommand{\NormalTok}[1]{#1}
\newcommand{\OperatorTok}[1]{\textcolor[rgb]{0.40,0.40,0.40}{#1}}
\newcommand{\OtherTok}[1]{\textcolor[rgb]{0.00,0.44,0.13}{#1}}
\newcommand{\PreprocessorTok}[1]{\textcolor[rgb]{0.74,0.48,0.00}{#1}}
\newcommand{\RegionMarkerTok}[1]{#1}
\newcommand{\SpecialCharTok}[1]{\textcolor[rgb]{0.25,0.44,0.63}{#1}}
\newcommand{\SpecialStringTok}[1]{\textcolor[rgb]{0.73,0.40,0.53}{#1}}
\newcommand{\StringTok}[1]{\textcolor[rgb]{0.25,0.44,0.63}{#1}}
\newcommand{\VariableTok}[1]{\textcolor[rgb]{0.10,0.09,0.49}{#1}}
\newcommand{\VerbatimStringTok}[1]{\textcolor[rgb]{0.25,0.44,0.63}{#1}}
\newcommand{\WarningTok}[1]{\textcolor[rgb]{0.38,0.63,0.69}{\textbf{\textit{#1}}}}
\newcommand{\passthrough}[1]{#1}

\begin{document}

\title{NTT-Friendly Pairing Curve Construction with\\Readability-Aware Prime Selection over a 1024-bit Field}

\author{
\IEEEauthorblockN{Tran Duy}
\IEEEauthorblockA{University of Engineering and Technology\\
Vietnam National University, Hanoi\\
Email: placeholder@vnu.edu.vn}
}

\maketitle

\begin{abstract}
Pairing-friendly elliptic curves are fundamental to advanced cryptographic protocols such as BLS signatures and zero-knowledge proof systems. However, the traditional Cocks-Pinch construction produces primes with no guaranteed algebraic structure, limiting their utility for NTT-based proof systems and hindering parameter auditability. We present three contributions: (1)~a modified Cocks-Pinch algorithm that constrains the cyclotomic seed to guarantee two-adicity of at least~36 in both the base field and scalar field primes, exceeding the industry-standard BLS12-381 curve; (2)~a multi-criteria readability scoring system that selects human-auditable primes from the valid candidate space without weakening security; and (3)~Curve1024, a concrete 1024-bit pairing-friendly curve ($E: y^2 = x^3 + 41$, $k=18$, 256-bit ECDLP security) with full parameter set and generator point published for public verification. A complete Rust implementation with 36~passing tests demonstrates practical signing operations and validates resistance against MOV, Anomalous, and TNFS attacks.
\end{abstract}

\begin{IEEEkeywords}
elliptic curve cryptography, pairing-friendly curves, Cocks-Pinch method, NTT-friendly primes, parameter auditability
\end{IEEEkeywords}

% Pandoc template for IEEE conference paper sections
\section{Introduction}

Pairing-friendly elliptic curves underpin a growing class of
cryptographic protocols, from BLS aggregate signatures \cite{BLS2001}
to succinct non-interactive zero-knowledge proofs (zk-SNARKs)
\cite{Groth2016, PLONK2019}. These protocols require curves whose
scalar field and base field support efficient polynomial arithmetic via
the Number Theoretic Transform (NTT), which in turn demands that both
field primes satisfy \(q \equiv 1 \pmod{2^s}\) for a sufficiently large
\(s\).

The Cocks-Pinch method \cite{CocksPinch2001} offers unmatched
flexibility in constructing pairing-friendly curves with arbitrary
embedding degrees. Unlike parametric families such as BN
\cite{FreemanScottTeske2010} or BLS12 \cite{BLS12_381}, the
Cocks-Pinch approach can target specific security levels and field
sizes. However, the traditional algorithm treats NTT-friendliness as an
afterthought---primes are generated with random two-adicity, typically
only 1--3 bits, rendering them unsuitable for modern proof systems.

A separate but related concern is \emph{parameter auditability}. The
Dual\_EC\_DRBG incident \cite{DualECDRBG} demonstrated that opaque
cryptographic constants can conceal backdoors. While the Cocks-Pinch
method is inherently transparent (all parameters derive from a public
seed), the resulting primes are typically inscrutable hexadecimal blobs
that resist manual inspection.

\textbf{Contributions.} This paper presents three results:

\begin{enumerate}
\def\labelenumi{\arabic{enumi}.}
\item
  \textbf{Improved Cocks-Pinch with NTT constraints.} We modify the
  Cocks-Pinch construction to \emph{guarantee} high two-adicity in both
  field primes by constraining the cyclotomic seed \(T\). Specifically,
  choosing \(T \equiv 0 \pmod{2^{12}}\) yields \(\nu_2(r-1) \geq 36\)
  deterministically (formalized in Lemma 1, Section 3), while an
  extended lift search ensures \(\nu_2(p-1) \geq 36\). In our
  construction, we obtain a two-adicity of 36, compared to BLS12-381's
  value of 32.
\item
  \textbf{Readability-aware prime selection.} We introduce a
  multi-criteria scoring system that evaluates candidate primes on eight
  structural properties (binary sparseness, zero limbs, hex density,
  etc.) and selects the most human-auditable parameters without
  sacrificing security.
\item
  \textbf{Curve1024: a concrete 1024-bit pairing-friendly curve.} We
  instantiate the above techniques to produce a complete, publicly
  verifiable curve with embedding degree \(k=18\), scalar field order
  \(r \approx 2^{512}\) (256-bit ECDLP security), and equation
  \(E: y^2 = x^3 + 41\). The full parameter set---including generator
  point---is published alongside a Rust implementation whose test suite
  verifies group arithmetic correctness and confirms resistance against
  MOV, Anomalous, and TNFS attacks.
\end{enumerate}

\textbf{Organization.} Section 2 reviews background on pairing-friendly
curves and the CM method. Section 3 presents our NTT-friendly
Cocks-Pinch modification. Section 4 describes the readability scoring
system. Section 5 presents the Curve1024 parameters and evaluates the
implementation. Section 6 discusses related work, and Section 7
concludes.

% Pandoc template for IEEE conference paper sections
\IEEEpubidadjcol
\section{Preliminaries}

\subsection{Elliptic Curves and
Pairings}\label{elliptic-curves-and-pairings}

Let \(\mathbb{F}_p\) be a prime field with \(p > 3\). An elliptic curve
in Short Weierstrass form is \(E: y^2 = x^3 + ax + b\) with
\(4a^3 + 27b^2 \neq 0\). The set \(E(\mathbb{F}_p)\) of rational points
forms an abelian group under the chord-and-tangent law
\cite{Silverman2009}.

A \emph{bilinear pairing} is a map
\(e: \mathbb{G}_1 \times \mathbb{G}_2 \to \mathbb{G}_T\) satisfying:

\begin{itemize}
\tightlist
\item
  \textbf{Bilinearity:} \(e(aP, bQ) = e(P, Q)^{ab}\) for all
  \(a, b \in \mathbb{Z}\)
\item
  \textbf{Non-degeneracy:} \(e(P, Q) \neq 1\) for generators
  \(P \in \mathbb{G}_1\), \(Q \in \mathbb{G}_2\)
\item
  \textbf{Computability:} \(e\) can be evaluated in polynomial time via
  Miller's algorithm \cite{Miller1986}
\end{itemize}

Here \(\mathbb{G}_1, \mathbb{G}_2\) are prime-order subgroups of
\(E(\mathbb{F}_p)\) and \(E(\mathbb{F}_{p^k})\) respectively, and
\(\mathbb{G}_T \subset \mathbb{F}_{p^k}^*\) is the target group. The
principal pairing constructions are the Weil pairing (symmetric,
computes two Miller loops) and the Tate pairing (one Miller loop plus a
final exponentiation). The \emph{optimal Ate pairing}
\cite{Vercauteren2010} further reduces the Miller loop length to
\(|T|\) bits (the curve seed), making it the preferred choice for
implementations.

\subsection{Embedding Degree and
Security}\label{embedding-degree-and-security}

A curve is \emph{pairing-friendly} with embedding degree \(k\) if there
exists a prime-order subgroup of size \(r\) such that \(r \mid p^k - 1\)
but \(r \nmid p^i - 1\) for \(0 < i < k\). The embedding degree
determines the pairing target field \(\mathbb{F}_{p^k}\): the MOV attack
\cite{MOV1993} reduces ECDLP in \(\mathbb{G}_1\) to DLP in
\(\mathbb{F}_{p^k}^*\), so \(k\) must be large enough that DLP in
\(\mathbb{F}_{p^k}\) remains intractable.

For a 1024-bit base field with \(k = 18\), the pairing target is
\(\mathbb{F}_{p^{18}}\)---a field of 18,432 bits. This far exceeds the
3,072-bit threshold for 128-bit DLP security and provides substantial
margin against the Tower Number Field Sieve (TNFS)
\cite{BarbulescuDuquesne2019, KimBarbulescu2016}. The choice \(k = 18\)
also admits the 18th cyclotomic polynomial
\(\Phi_{18}(T) = T^6 - T^3 + 1\), which produces 512-bit primes \(r\)
from 86-bit seeds---an efficient ratio for the Cocks-Pinch construction.

\subsection{Complex Multiplication
Method}\label{complex-multiplication-method}

The CM method constructs curves with a prescribed group order. Given a
CM discriminant \(D < 0\), the existence of a curve over
\(\mathbb{F}_p\) with Frobenius trace \(t\) requires the Diophantine
equation:

\[4p = t^2 + |D| \cdot y^2\]

to have an integer solution. The \(j\)-invariant is obtained from the
Hilbert class polynomial \(H_D(x)\). For \(D = -3\), we have
\(H_{-3}(x) = x\), giving \(j = 0\) and the simple curve form
\(y^2 = x^3 + b\).

\subsection{Traditional Cocks-Pinch
Algorithm}\label{traditional-cocks-pinch-algorithm}

The Cocks-Pinch method \cite{CocksPinch2001} constructs
pairing-friendly curves by starting from the subgroup order \(r\) rather
than the field prime \(p\). For embedding degree \(k = 18\):

\begin{enumerate}
\def\labelenumi{\arabic{enumi}.}
\tightlist
\item
  \textbf{Generate \(r\)}: Choose a parameter \(T\) and compute
  \(r = \Phi_{18}(T) = T^6 - T^3 + 1\). Verify that \(r\) is prime.
\item
  \textbf{Compute square root}: Find \(\beta = \sqrt{-3} \pmod{r}\) via
  Tonelli-Shanks.
\item
  \textbf{CM parameters}: For each \(i\) with \(\gcd(i, 18) = 1\),
  compute \(t_0 \equiv T^i + 1 \pmod{r}\) and
  \(y_0 \equiv (t_0 - 2)/\beta \pmod{r}\).
\item
  \textbf{Lift to integers}: Search over \((h_t, h_y) \in [-20, 20]^2\)
  for \(t = t_0 + h_t \cdot r\), \(y = y_0 + h_y \cdot r\) such that
  \(p = (t^2 + 3y^2)/4\) is prime of the target size.
\end{enumerate}

The resulting curve has \(\rho = \log p / \log r \approx 2\), which is
the standard tradeoff of the Cocks-Pinch method compared to parametric
families (\(\rho = 1\) for BN, \(\rho = 1.5\) for BLS12).

\subsection{Two-Adicity and NTT}\label{two-adicity-and-ntt}

The \emph{two-adicity} of a prime \(q\) is \(\nu_2(q-1)\), the largest
power of 2 dividing \(q-1\). A field \(\mathbb{F}_q\) with two-adicity
\(s\) contains a primitive \(2^s\)-th root of unity, enabling NTT of
length up to \(2^s\). Modern ZK proof systems (Groth16
\cite{Groth2016}, PLONK \cite{PLONK2019}) require NTT for polynomial
multiplication over both \(\mathbb{F}_r\) (constraint evaluation) and
\(\mathbb{F}_p\) (pairing computation). BLS12-381 achieves two-adicity
32 in both fields \cite{BLS12_381}.

% Pandoc template for IEEE conference paper sections
\section{Improved Cocks-Pinch with NTT Constraints}

We modify the traditional Cocks-Pinch algorithm to guarantee high
two-adicity in both output primes. The key insight is that the algebraic
structure of the cyclotomic polynomial \(\Phi_{18}\) allows
deterministic control over the two-adicity of \(r\), while an extended
lift search achieves the same for \(p\).

\subsection{NTT-Friendly Scalar Field via Seed
Constraint}\label{ntt-friendly-scalar-field-via-seed-constraint}

Since \(r = \Phi_{18}(T) = T^6 - T^3 + 1\), we observe:

\[r - 1 = T^6 - T^3 = T^3(T^3 - 1)\]

If \(T \equiv 0 \pmod{2^m}\), then \(T^3 \equiv 0 \pmod{2^{3m}}\). Since
\(T^3 - 1\) is odd (as \(T^3\) is even), we obtain:

\[\nu_2(r - 1) = \nu_2(T^3) = 3m\]

\textbf{Consequence:} To guarantee \(\nu_2(r-1) \geq s\), it suffices to
choose \(T \equiv 0 \pmod{2^{\lceil s/3 \rceil}}\). This is a
\emph{deterministic} guarantee requiring no rejection sampling.

For our target of \(s = 36\), we set \(m = \lceil 36/3 \rceil = 12\),
constraining \(T \equiv 0 \pmod{2^{12}}\). The search space remains
\(2^{86-12} = 2^{74}\) candidates---more than sufficient for finding
prime \(r\).

\subsection{NTT-Friendly Base Field via Extended Lift
Search}\label{ntt-friendly-base-field-via-extended-lift-search}

After obtaining a valid \(r\) with the desired two-adicity, the base
field prime is:

\[p = \frac{(t_0 + h_t r)^2 + 3(y_0 + h_y r)^2}{4}\]

The two-adicity of \(p-1\) depends on the lift parameters
\((h_t, h_y)\). To ensure \(\nu_2(p-1) \geq s_p\), we extend the search
grid beyond the default \([-20, 20]^2\):

\[\text{half\_range} = 20 + 2^{\max(0,\, s_p - 3m)}\]

This extension guarantees sufficient coverage of bit patterns in \(p-1\)
at positions above \(3m\). For \(s_p = 36\) and \(m = 12\) (so
\(3m = 36\)), the grid remains at \([-21, 21]^2\)---only marginally
larger than the default.

\subsection{Algorithm}\label{algorithm}

The complete NTT-friendly Cocks-Pinch construction proceeds as follows:

\textbf{Input:} Embedding degree \(k=18\), target sizes \((|r|, |p|)\),
minimum two-adicity \(s\)

\textbf{Output:} Primes \((p, r)\) with \(\nu_2(r-1) \geq s\) and
\(\nu_2(p-1) \geq s\)

\begin{enumerate}
\def\labelenumi{\arabic{enumi}.}
\tightlist
\item
  Set \(m \leftarrow \lceil s/3 \rceil\) and
  \(\text{step} \leftarrow 2^m\)
\item
  Repeat:

  \begin{enumerate}
  \def\labelenumii{\alph{enumii}.}
  \tightlist
  \item
    Choose \(T\) as a random multiple of step in the target range
  \item
    Compute \(r \leftarrow T^6 - T^3 + 1\)
  \item
    If \(r\) is not prime or \(|r| \neq\) target, continue
  \item
    Compute \(\beta \leftarrow \sqrt{-3} \pmod{r}\) (Tonelli-Shanks)
  \item
    For each \(i\) with \(\gcd(i, 18) = 1\):

    \begin{itemize}
    \tightlist
    \item
      \(t_0 \leftarrow T^i + 1 \pmod{r}\)
    \item
      \(y_0 \leftarrow (t_0 - 2) \cdot \beta^{-1} \pmod{r}\)
    \item
      \(\text{range} \leftarrow 20 + 2^{\max(0, s - 3m)}\)
    \item
      For \((h_t, h_y) \in [-\text{range}, \text{range}]^2\):

      \begin{itemize}
      \tightlist
      \item
        \(t \leftarrow t_0 + h_t \cdot r\);
        \(y \leftarrow y_0 + h_y \cdot r\)
      \item
        \(p \leftarrow (t^2 + 3y^2) / 4\)
      \item
        If \(|p| =\) target AND \(\nu_2(p-1) \geq s\) AND \(p\) is
        prime: return \((p, r, t, y)\)
      \end{itemize}
    \end{itemize}
  \end{enumerate}
\item
  Until max attempts reached
\end{enumerate}

\subsection{Complexity Analysis}\label{complexity-analysis}

The following table compares the traditional and improved algorithms.
The NTT constraint increases the average number of attempts by
approximately 10--60\(\times\) (from \textasciitilde100--500 to
\textasciitilde3,000--6,000), but the absolute generation time remains
under one minute---a one-time cost amortized over the lifetime of the
curve parameters.

\begin{table}[t]
\centering
\caption{Traditional vs.~Improved Cocks-Pinch}
\begin{tabular}{lll}
\toprule
Metric & Traditional & Improved \\
\midrule
Constraint on \(T\) & None & \(T \equiv 0 \pmod{2^{12}}\) \\
\(\nu_2(r-1)\) & Random (1--3) & \(\geq 36\) (guaranteed) \\
\(\nu_2(p-1)\) & Random (1--3) & \(\geq 36\) (with ext. grid) \\
Grid size & \([-20,20]^2\) & \([-21,21]^2\) \\
Avg. attempts & 100--500 & 3,000--6,000 \\
Generation time & 2--10 s & 30--60 s \\
NTT support & No & Up to \(2^{36}\) \\
\bottomrule
\end{tabular}
\end{table}

\subsection{Comparison with BLS12-381}\label{comparison-with-bls12-381}

BLS12-381 \cite{BLS12_381}, the current industry standard for ZK proof
systems, achieves two-adicity 32 in both fields. Our construction
achieves 36, enabling NTT of length
\(2^{36} \approx 6.8 \times 10^{10}\)---a \(16\times\) improvement in
maximum polynomial degree. The tradeoff is a larger \(\rho\)-value (2.0
vs.~1.5) and correspondingly larger field elements, which is acceptable
for applications prioritizing long-term security (256-bit ECDLP
vs.~128-bit for BLS12-381).

\input{chapter/04-readability}
% Pandoc template for IEEE conference paper sections
\section{Curve1024: Parameters and Evaluation}

This section presents the concrete Curve1024 parameters produced by our
improved Cocks-Pinch construction with readability-aware selection,
followed by implementation details and security analysis.

\subsection{Curve1024 Parameters}\label{curve1024-parameters}

The construction produces the following pairing-friendly elliptic curve
over a 1024-bit prime field. All parameters are publicly verifiable from
the seed \(T\).

\subsubsection{Seed and Construction}\label{seed-and-construction}

The curve is fully determined by a single 86-bit seed:

\[T = \texttt{0x26704d2ace0facdd539} \cdot 2^{12}\]

This seed satisfies \(T \equiv 0 \pmod{2^{12}}\), guaranteeing
\(\nu_2(r-1) = 36\) via the factorization \(r - 1 = T^3(T^3 - 1)\).

\subsubsection{Scalar Field Order (512
bits)}\label{scalar-field-order-512-bits}

The scalar field prime \(r = \Phi_{18}(T) = T^6 - T^3 + 1\) decomposes
as:

\[r = d_r \cdot 2^{36} + 1\]

where \(d_r\) is a 476-bit odd integer. In 64-bit limb representation
(MSB first):

\begin{table}[t]
\centering
\caption{Scalar field prime \(r\) --- limb decomposition}
\begin{tabular}{lll}
\toprule
Limb & Value & Note \\
\midrule
15 & \texttt{0000000000000000} & zero \\
14 & \texttt{0000000000000000} & zero \\
13 & \texttt{0000000000000000} & zero \\
12 & \texttt{0000000000000000} & zero \\
11 & \texttt{0000000000000000} & zero \\
10 & \texttt{0000000000000000} & zero \\
9 & \texttt{0000000000000000} & zero \\
8 & \texttt{0000000000000000} & zero \\
7 & \texttt{c042c3389e72044d} &  \\
6 & \texttt{9ec8078aea7bc954} &  \\
5 & \texttt{ebf209d918f0ef27} &  \\
4 & \texttt{425259e247711927} &  \\
3 & \texttt{97463bf23832d0de} &  \\
2 & \texttt{5e46e47fc61cab50} &  \\
1 & \texttt{4bcfd36fc2a945c5} &  \\
0 & \texttt{1055d97000000001} & tail: \(d \cdot 2^{36} + 1\) \\
\bottomrule
\end{tabular}
\end{table}

\textbf{Key readability features:} 8 of 16 limbs are zero (50\%), and
the final limb visually confirms the NTT-friendly structure with the
\texttt{00000001} suffix (36 trailing zero bits before the \(+1\)).

\subsubsection{Base Field Prime (1024
bits)}\label{base-field-prime-1024-bits}

The base field prime \(p = (t^2 + 3y^2)/4\) decomposes as:

\[p = d_p \cdot 2^{36} + 1\]

where \(d_p\) is a 988-bit odd integer. The final limb shares the same
tail pattern as \(r\):

\begin{table}[t]
\centering
\caption{Base field prime \(p\) --- limb decomposition}
\begin{tabular}{lll}
\toprule
Limb & Value & Note \\
\midrule
15 & \texttt{c0859da83a500988} &  \\
14 & \texttt{10eea35c61fea054} &  \\
13 & \texttt{c866902599220d30} &  \\
12 & \texttt{c93f38265ed6b830} &  \\
11 & \texttt{5785b8a073e482e8} &  \\
10 & \texttt{685cfeeb388dc45e} &  \\
9 & \texttt{bca1686245b326fe} &  \\
8 & \texttt{824ca74ac0844160} &  \\
7 & \texttt{f4e289806084e50e} &  \\
6 & \texttt{6c02f32cb94d786b} &  \\
5 & \texttt{f80fb080c07d6a07} &  \\
4 & \texttt{73f2b4591eb70df3} &  \\
3 & \texttt{9367aa32ebb42449} &  \\
2 & \texttt{e0a648f61a20e876} &  \\
1 & \texttt{f971f3bf976a7d6a} &  \\
0 & \texttt{dd20d97000000001} & tail: \(d \cdot 2^{36} + 1\) \\
\bottomrule
\end{tabular}
\end{table}

\textbf{Shared tail pattern:} Both \(p\) and \(r\) end with
\texttt{...d970\_00000001}, a direct consequence of the seed constraint
propagating through the construction. This visual consistency aids
manual verification.

\subsubsection{Curve Summary}\label{curve-summary}

\begin{table}[t]
\centering
\caption{Curve1024 parameters}
\begin{tabular}{ll}
\toprule
Parameter & Value \\
\midrule
Equation & \(E: y^2 = x^3 + 41\) \\
Base field \(|p|\) & 1024 bits \\
Scalar field \(|r|\) & 512 bits \\
Embedding degree \(k\) & 18 \\
CM discriminant \(D\) & \(-3\) (hence \(j = 0\)) \\
\(\rho = \log p / \log r\) & 2.0 \\
\(\nu_2(p-1)\) & 36 \\
\(\nu_2(r-1)\) & 36 \\
Cofactor \(h = \#E/r\) & \(\approx 2^{512}\) \\
\bottomrule
\end{tabular}
\end{table}

The curve equation \(y^2 = x^3 + 41\) arises from the CM method with
\(D = -3\): the Hilbert class polynomial \(H_{-3}(x) = x\) gives
\(j = 0\), yielding the family \(y^2 = x^3 + b\). The twist \(b = 41\)
is selected by verifying which of the six sextic twists has group order
divisible by \(r\).

A generator point \(G = (G_x, G_y)\) of prime order \(r\) is obtained by
cofactor multiplication of a random curve point. The full coordinates
are published in the project repository \cite{LumenMath}.

\subsubsection{Readability Metrics}\label{readability-metrics}

\begin{table}[t]
\centering
\caption{Readability comparison}
\begin{tabular}{lll}
\toprule
Metric & \(p\) (1024-bit) & \(r\) (512-bit) \\
\midrule
Hamming weight & 464/1024 (45\%) & 235/512 (46\%) \\
Two-adicity & 36 & 36 \\
Zero 64-bit limbs & 0/16 & 8/16 (50\%) \\
Hex zero digits & 33/256 (13\%) & 16/128 (12.5\%) \\
Trailing zero bits & 36 & 36 \\
\bottomrule
\end{tabular}
\end{table}

The scalar field prime \(r\) is notably sparse: half its limb
representation is zero, making it easy to recognize and verify in
debugging output. Both primes exceed BLS12-381's two-adicity of 32.

\subsection{Implementation}\label{implementation}

We implement the complete system in Rust without external cryptographic
dependencies \cite{LumenMath}. The implementation comprises:

\begin{itemize}
\tightlist
\item
  A custom \texttt{U1024} big integer type with Montgomery
  multiplication \cite{Montgomery1985}
\item
  Prime field arithmetic (\texttt{PrimeFieldElement}) with constant-time
  Fermat inversion
\item
  Elliptic curve group operations in affine coordinates
\item
  Schnorr and ECDSA signature schemes
\item
  The parameter generation tool (improved Cocks-Pinch + readability
  scoring)
\end{itemize}

The test suite contains 36 tests covering arithmetic correctness, group
law verification, and attack simulation (MOV, Anomalous, Invalid Curve).

\subsection{Performance Benchmarks}\label{performance-benchmarks}

\begin{table}[t]
\centering
\caption{Performance benchmarks (Criterion, release build, LLVM O3, no AVX-512)}
\begin{tabular}{llll}
\toprule
Operation & Min & Mean & Max \\
\midrule
Key generation & 1.629 s & 1.687 s & 1.757 s \\
Schnorr sign & 1.473 s & 1.551 s & 1.641 s \\
Schnorr verify & 2.656 s & 2.831 s & 3.025 s \\
ECDSA sign & 1.524 s & 1.606 s & 1.711 s \\
ECDSA verify & 2.847 s & 3.140 s & 3.459 s \\
Scalar mult. \([k]G\) & 1.061 s & \textasciitilde1.2 s & 1.489 s \\
\bottomrule
\end{tabular}
\end{table}

The performance reflects the 1024-bit field size: each scalar
multiplication requires \textasciitilde512 double-and-add iterations,
each involving a Fermat inversion (\textasciitilde1536 Montgomery
multiplications). This is approximately \(8000\times\) slower than
optimized secp256k1 implementations, which is expected given the
\(4\times\) larger field and \(O(n^2)\) multiplication cost. The system
targets offline signing scenarios (legal documents, software releases,
PKI certificates) where seconds of latency are acceptable.

\subsection{Security Analysis}\label{security-analysis}

\begin{table}[t]
\centering
\caption{Security evaluation summary}
\begin{tabular}{lll}
\toprule
Attack & Complexity & Status \\
\midrule
Pollard-\(\rho\) (ECDLP) & \(O(2^{256})\) & Secure \\
MOV reduction & DLP in \(\mathbb{F}_{p^{18}}\) (18432-bit) & Secure \\
Anomalous (SSSA) & Requires \(\#E = p\) & Immune \\
TNFS {[}@BarbulescuDuquesne2019{]} & \(\geq 250\)-bit in \(\mathbb{F}_{p^{18}}\) & Secure \\
Invalid curve & Point validation in constructor & Mitigated \\
\bottomrule
\end{tabular}
\end{table}

The 256-bit ECDLP security exceeds NIST SP 800-57 recommendations
through 2030+ \cite{NIST80057}. The MOV attack \cite{MOV1993} maps
ECDLP to DLP in \(\mathbb{F}_{p^{18}}\), a field of 18,432 bits---far
exceeding the 3,072-bit threshold for 128-bit security. The curve is
provably non-anomalous \cite{SmartAnomaly1999} since
\(\#E(\mathbb{F}_p) = h \cdot r\) with
\(r \approx 2^{512} \neq p \approx 2^{1024}\).

\textbf{Limitation:} The current affine-coordinate implementation uses
variable-time double-and-add, which leaks scalar bits through timing. A
Montgomery ladder (constant-time) is identified as future work.

\subsection{Curve Family Comparison}\label{curve-family-comparison}

\begin{table}[t]
\centering
\caption{Comparison with standard pairing-friendly curve families}
\begin{tabular}{llllll}
\toprule
Curve & \(|p|\) (bits) & \(\rho\) & \(\nu_2\) & ECDLP Security & Pairing Target \\
\midrule
BN254 & 254 & 1.0 & \textasciitilde1 & 100-bit & 3048-bit \\
BLS12-381 & 381 & 1.5 & 32 & 128-bit & 4572-bit \\
KSS18 (generic) & varies & 1.33 & random & varies & varies \\
\textbf{Curve1024} & \textbf{1024} & 2.0 & \textbf{36} & \textbf{256-bit} & \textbf{18432-bit} \\
\bottomrule
\end{tabular}
\end{table}

Curve1024 occupies a unique position: it provides the highest ECDLP
security level (256-bit) and the highest two-adicity (36) among deployed
pairing-friendly curves, at the cost of a larger \(\rho\)-value inherent
to the Cocks-Pinch method. The 18,432-bit pairing target field provides
substantial margin against TNFS attacks.

% Pandoc template for IEEE conference paper sections
\section{Related Work}

\textbf{Pairing-friendly curve taxonomy.} Freeman, Scott, and Teske
\cite{FreemanScottTeske2010} provide a comprehensive classification of
pairing-friendly curve constructions, including the Cocks-Pinch method,
MNT curves, and polynomial families (BN, BLS, KSS). Their work
establishes the theoretical framework but does not address
NTT-friendliness as a design criterion.

\textbf{BN curves and TNFS downgrade.} Barreto-Naehrig curves
\cite{BarretoNaehrig2005} achieve \(\rho = 1\) with embedding degree
12, making BN254 the standard for Ethereum's precompiles. However, Kim
and Barbulescu \cite{KimBarbulescu2016} demonstrated that the extended
Tower Number Field Sieve reduces BN254's effective security from 128-bit
to approximately 100-bit. This TNFS vulnerability motivates the move to
larger fields: our 1024-bit base field with \(k=18\) yields an
18,432-bit pairing target, providing substantial margin against all
known NFS variants.

\textbf{TNFS security updates.} Barbulescu and Duquesne
\cite{BarbulescuDuquesne2019} provided updated key size estimations
showing that many deployed pairing-friendly curves offer less security
than originally claimed. Their analysis confirms that embedding degree
18 with a 1024-bit prime provides \(\geq 250\)-bit security in the
pairing target field.

\textbf{BLS12-381.} Bowe \cite{BLS12_381} designed BLS12-381
specifically for zk-SNARKs, achieving two-adicity 32 in both fields
through careful parameter selection within the BLS12 family. Our work
achieves higher two-adicity (36) using the more flexible Cocks-Pinch
method, at the cost of a larger \(\rho\)-value (2.0 vs.~1.5).

\textbf{KSS-18 curves.} Kachisa, Schaefer, and Scott \cite{KSS2008}
construct pairing-friendly curves with embedding degree 18 using the
Brezing-Weng method \cite{BrezingWeng2005}, achieving
\(\rho \approx 1.33\). While KSS-18 offers better bandwidth efficiency
than Cocks-Pinch (\(\rho = 2\)), the parametric family does not
inherently guarantee NTT-friendliness, and the fixed polynomial
structure limits the ability to optimize for readability across a large
candidate space.

\textbf{Pasta curves.} The Pallas/Vesta cycle \cite{PastaCurves}
achieves NTT-friendliness through a 2-cycle of curves, enabling
recursive proof composition. Unlike our approach, Pasta curves are not
pairing-friendly and target a different application domain (Halo 2
recursive proofs).

\textbf{Nothing-up-my-sleeve numbers and SafeCurves.} The cryptographic
community has long advocated for transparent parameter generation to
demonstrate absence of backdoors. The SafeCurves project
\cite{SafeCurves} formalizes criteria including rigidity, twist
security, and completeness. Our readability scoring takes a
complementary approach: rather than deriving primes from a fixed public
constant, we select the most inspectable prime from a large space of
equally secure candidates, while the Cocks-Pinch seed provides full
reproducibility.

\textbf{Cocks-Pinch extensions.} Chen, Hong, and Wang
\cite{CocksPinchK5K8} extend the Cocks-Pinch method to embedding
degrees 5--8 with optimal Ate pairing computation. Our work focuses on
\(k=18\) with the novel addition of NTT constraints and readability
scoring.

% Pandoc template for IEEE conference paper sections
\section{Conclusion}

We presented three contributions to pairing-friendly elliptic curve
construction:

\begin{enumerate}
\def\labelenumi{\arabic{enumi}.}
\item
  \textbf{NTT-friendly Cocks-Pinch.} By constraining the cyclotomic seed
  \(T \equiv 0 \pmod{2^{12}}\), we guarantee two-adicity \(\geq 36\) in
  the scalar field deterministically, while an extended lift search
  achieves the same for the base field. This exceeds the
  industry-standard BLS12-381 (two-adicity 32) at a modest increase in
  generation time (30--60 s vs.~2--10 s).
\item
  \textbf{Readability-aware prime selection.} Our multi-criteria scoring
  system produces primes that are more amenable to human inspection
  without any security degradation, addressing the auditability concern
  in cryptographic parameter generation.
\item
  \textbf{Curve1024.} We instantiate these techniques into a concrete,
  publicly verifiable curve: \(E: y^2 = x^3 + 41\) over a 1024-bit prime
  field with embedding degree 18, providing 256-bit ECDLP security and
  NTT support up to length \(2^{36}\) in both fields. The complete
  parameter set (field primes, generator point, CM data) is published
  and validated by a Rust implementation with 36 passing tests.
\end{enumerate}

The system is suitable for offline signing applications requiring
long-term security guarantees and for future pairing-based protocols
(BLS aggregate signatures, zk-SNARKs) that benefit from the NTT-friendly
field structure.

\textbf{Future work} includes: (1) projective/Jacobian coordinates to
eliminate per-operation inversions, (2) constant-time Montgomery ladder
for side-channel resistance, (3) full optimal Ate pairing implementation
over \(\mathbb{F}_{p^{18}}\), and (4) extension to other embedding
degrees (\(k = 12, 24\)).

\subsection{Availability}\label{availability}

The implementation and all curve parameters are available as open-source
software under the MIT license \cite{LumenMath}.


\bibliographystyle{IEEEtran}
\bibliography{references}

\end{document}
